\documentclass[a4paper,12pt]{article}
\usepackage{iftex}
\ifPDFTeX
  \usepackage[T2A]{fontenc}
  \usepackage[utf8]{inputenc}
  \usepackage[russian]{babel}
  \renewcommand{\familydefault}{\sfdefault}
\else
  \usepackage{fontspec}
  \setmainfont{Inter}
  \usepackage[russian]{babel}
\fi
\usepackage{amsmath,amssymb,mathtools}
\usepackage{icomma}
\usepackage{geometry}
\geometry{left=21mm,right=21mm,top=21mm,bottom=23mm}
\usepackage{microtype}
\usepackage{tikz}
\usetikzlibrary{arrows.meta,positioning,shapes.geometric,calc}
\usepackage{tabularx,booktabs,array}
\usepackage{enumitem}
\usepackage{hyperref}
\usepackage{fancyhdr}
\usepackage{setspace}
\usepackage{xcolor}
\usepackage{needspace}
\definecolor{accent}{HTML}{007A7A}
\definecolor{darkaccent}{HTML}{004E57}
\definecolor{ink}{HTML}{252F34}
\definecolor{muted}{HTML}{5D686E}
\hypersetup{hidelinks,pdftitle={Чек-лист. Деревья вероятностей и сложные текстовые задачи},pdfauthor={Лёвин Артём Александрович}}
\setstretch{1.14}
\setlength{\parindent}{0pt}
\setlength{\parskip}{5pt}
\emergencystretch=3em
\setlist[itemize]{topsep=4pt,itemsep=3pt,leftmargin=17pt}
\setlist[enumerate]{topsep=4pt,itemsep=4pt,leftmargin=19pt}
\pagestyle{fancy}
\fancyhf{}
\fancyhead[L]{\small\textcolor{muted}{Теория вероятностей}}
\fancyhead[R]{\small\textcolor{muted}{Екатерина Гнедкова}}
\fancyfoot[L]{\footnotesize\textcolor{muted}{Артём Александрович Лёвин}}
\fancyfoot[R]{\footnotesize\textcolor{muted}{\thepage}}
\renewcommand{\headrulewidth}{0.3pt}
\renewcommand{\footrulewidth}{0pt}
\setlength{\headheight}{16pt}
\newcommand{\smalltopic}[1]{\par\vspace{7pt}\noindent\textcolor{darkaccent}{\bfseries #1}\par\vspace{1pt}}
\newcommand{\important}[1]{\textcolor{darkaccent}{\bfseries #1}}
\newcommand{\ex}[1]{\needspace{9\baselineskip}\smalltopic{#1}}
\newcommand{\answer}[2]{#1 & #2\\}
\begin{document}
\begin{titlepage}
\thispagestyle{empty}
\begin{center}
\vspace*{26mm}
{\large\bfseries\textcolor{darkaccent}{МАТЕМАТИКА}}\\[19mm]
{\fontsize{30}{36}\selectfont\bfseries Чек-лист}\\[8mm]
{\fontsize{18}{24}\selectfont Деревья вероятностей\\[-1mm]и сложные текстовые задачи}\\[13mm]
{\large Подготовка к ЕГЭ по профильной математике}\\[29mm]
{\Large\bfseries Екатерина Гнедкова}\\[12mm]
{\large 17 ноября 2026 г.}
\vfill
{\normalsize Лёвин Артём Александрович}\\[1.5mm]
{\small эксклюзивно для Екатерины Гнедковой}\\[13mm]
\begin{tikzpicture}
\draw[accent!70,line width=0.85pt] (-4,0) .. controls (-1.75,0.65) and (1.55,-0.72) .. (4,0.13);
\end{tikzpicture}
\end{center}
\end{titlepage}

\section*{Главная идея занятия}
В сложной задаче на вероятность важно сначала разобраться, \emph{какие случайные события происходят и в каком порядке}. После этого длинное условие превращается в дерево исходов, таблицу или уравнение с неизвестной долей.

\smalltopic{Что необходимо уметь после повторения}
\begin{itemize}
\item Выделять последовательные этапы опыта и строить дерево вероятностей.
\item Находить вероятность конкретного пути, а затем объединять подходящие пути.
\item Отличать действительное состояние объекта от результата проверки или наблюдения.
\item Составлять уравнения по общей вероятности события и находить неизвестные доли.
\item Выбирать таблицу исходов для небольшого числа игр или испытаний.
\item Решать задачи на повторные попытки и на выбор предметов без возвращения.
\end{itemize}

\section{Три правила, на которых держится решение}
\smalltopic{1. Дополнительное событие}
Если событие $A$ происходит с вероятностью $p$, то противоположное событие $\overline A$ происходит с вероятностью
\[
P(\overline A)=1-p.
\]
Например, если в одной попытке вероятность успеха $0{,}8$, вероятность неудачи равна $0{,}2$.

\smalltopic{2. Последовательность событий: умножаем}
Если требуется, чтобы последовательно произошли $A$ и затем $B$, используем
\[
P(A\cap B)=P(A)\,P(B\mid A).
\]
Здесь $P(B\mid A)$ означает вероятность события $B$ при условии, что событие $A$ уже произошло. Если события независимы, она совпадает с $P(B)$.

\smalltopic{3. Альтернативные непересекающиеся пути: складываем}
Если искомый результат достигается несколькими \emph{взаимоисключающими} способами, то
\[
P(A_1\cup A_2\cup\dots)=P(A_1)+P(A_2)+\dots.
\]
Например, цель уничтожена \emph{либо} первым выстрелом, \emph{либо} вторым после первого промаха. Одновременно эти два исхода в модели с прекращением стрельбы после успеха не происходят.

\important{Памятка:} вдоль одной траектории дерева вероятности перемножаются; вероятности разных подходящих непересекающихся траекторий складываются. В каждой развилке сумма вероятностей всех выходящих ветвей равна~$1$.

\section{Как перевести условие в дерево}
\begin{enumerate}
\item Назовите этапы случайного опыта. Например: \emph{действительно болен / здоров} $\to$ \emph{анализ положителен / отрицателен}.
\item Нарисуйте развилки только для действительно возможных продолжений опыта. Если после положительного анализа повторный тест не проводится, соответствующую ветку продолжать не нужно.
\item Подпишите вероятность на \emph{каждой} ветви. Учитывайте, меняется ли она после предыдущего действия.
\item Отметьте все пути, которые приводят к требуемому результату.
\item Перемножьте вероятности вдоль каждого выделенного пути и сложите результаты для непересекающихся путей.
\end{enumerate}

\ex{Пример 1. Повторная диагностика здоровой собаки}
У здоровой собаки один анализ ошибочно оказывается положительным с вероятностью $0{,}005$. Если анализ отрицательный, проводится ещё один. При повторном отрицательном результате исследование заканчивается. Какова вероятность, что хотя бы один из двух анализов окажется положительным?

\begin{center}
\begin{tikzpicture}[x=1cm,y=1cm,>=Latex,every node/.style={font=\small}]
\node (s) at (0,0) {Первый анализ};
\node (p) at (-3,-1.4) {$+$}; \node (n) at (3,-1.4) {$-$};
\node (p2) at (1.55,-2.95) {$+$}; \node (n2) at (4.75,-2.95) {$-$};
\draw[->,thick,accent] (s) -- node[above left] {$0{,}005$} (p);
\draw[->,thick,accent] (s) -- node[above right] {$0{,}995$} (n);
\draw[->,thick,accent] (n) -- node[above left] {$0{,}005$} (p2);
\draw[->,thick,accent] (n) -- node[above right] {$0{,}995$} (n2);
\node[below=2mm of p] {Диагноз $+$};
\node[below=2mm of p2] {Диагноз $+$};
\end{tikzpicture}
\end{center}
Возможны два подходящих пути: положительный первый анализ \emph{или} отрицательный первый и положительный второй:
\[
P(+)=0{,}005+0{,}995\cdot0{,}005=0{,}009975\approx0{,}01.
\]
\important{Точность формулировки:} слово «ошибочно» имеет смысл: собака изначально здорова, поэтому положительный анализ является ложноположительным. В дереве мы отмечаем именно результат «$+$».

\ex{Пример 2. Два дня и два возможных состояния погоды}
Погода бывает хорошей и замечательной. С вероятностью $0{,}8$ погода завтра останется такой же, как сегодня; с вероятностью $0{,}2$ изменится. Сегодня замечательная погода. Найдём вероятность замечательной погоды послезавтра.

\[
P(\text{замечательная через два дня})=
\underbrace{0{,}8\cdot0{,}8}_{\text{без изменений}}+
\underbrace{0{,}2\cdot0{,}2}_{\text{две смены}}=0{,}68.
\]
\important{Частая ошибка:} учитывайте оба пути к конечному состоянию. Вернуться к исходной погоде можно и после двух последовательных смен.

\section{Результат анализа и действительное состояние}
При медицинской диагностике полезно разделять \emph{два уровня}: действительное состояние человека и вывод прибора. Положительный анализ возможен и у больного, и у здорового.

Обозначим: $B$~--- человек действительно болен, $Z$~--- здоров, $+$~--- положительный результат. Тогда
\[
\boxed{P(+)=P(B)P(+\mid B)+P(Z)P(+\mid Z).}
\]
Это \important{формула полной вероятности}: она складывает два непересекающихся способа получить положительный результат.

\ex{Пример 3. Известна доля действительно больных}
Среди пациентов с подозрением на гепатит действительно больны $5\%$. У больного пациента анализ положителен с вероятностью $0{,}9$, у здорового~--- с вероятностью $0{,}01$. Какова вероятность положительного анализа у случайного пациента?
\[
P(+)=0{,}05\cdot0{,}9+0{,}95\cdot0{,}01
=0{,}045+0{,}0095=\boxed{0{,}0545}.
\]

\ex{Пример 4. Известна общая доля положительных анализов}
Вероятности положительного результата у больного и здорового те же: $0{,}9$ и $0{,}01$. Но теперь известно, что положительный результат получают $6\%$ всех пациентов. Нужно выяснить долю действительно больных.

Пусть $x=P(B)$, тогда $P(Z)=1-x$. Общая вероятность положительного результата известна:
\[
0{,}9x+0{,}01(1-x)=0{,}06.
\]
Раскрываем скобки и решаем уравнение:
\[
0{,}89x=0{,}05,\qquad x=\frac{5}{89}\approx0{,}05618.
\]
Итак, действительно болеют приблизительно $5{,}62\%$ пациентов.

\important{Ключевое различие:} «$6\%$ пациентов имеют положительный анализ» описывает \emph{наблюдаемый результат}. Эта фраза сама по себе не означает, что действительно больны $6\%$ пациентов.

\section{Обратная задача о доле группы}
Тот же приём применим, когда предметы поступают с разных заводов или хозяйств, а люди относятся к разным демографическим группам.

Пусть доля первой группы~--- $x$, второй~--- $1-x$. Если событие внутри первой группы происходит с вероятностью $a$, а внутри второй~--- с вероятностью $b$, то общая вероятность равна
\[
q=ax+b(1-x).
\]
При $a\ne b$ доля первой группы находится по формуле
\[
\boxed{x=\frac{q-b}{a-b}}.
\]
Обязательно проверьте, что найденное значение лежит между $0$ и $1$.

\ex{Пример 5. Яйца из двух хозяйств}
Агрофирма закупает яйца в двух хозяйствах. В первом $40\%$ яиц высшей категории, во втором $20\%$. Среди всех купленных яиц высшая категория составляет $35\%$. Какова доля яиц из первого хозяйства?
\[
0{,}40x+0{,}20(1-x)=0{,}35
\quad\Longrightarrow\quad 0{,}20x=0{,}15
\quad\Longrightarrow\quad \boxed{x=0{,}75}.
\]
Следовательно, $75\%$ всех яиц закуплено у первого хозяйства.

\ex{Пример 6. Изделия двух заводов}
Первый завод выпускает $2\%$ бракованных ламп, второй~--- $3\%$. Среди случайно выбранных ламп этого типа брак составляет $2{,}4\%$. Если $x$~--- доля изделий первого завода, то
\[
0{,}02x+0{,}03(1-x)=0{,}024
\quad\Longrightarrow\quad \boxed{x=0{,}60}.
\]
Значит, вероятность, что случайная лампа изготовлена первым заводом, равна $0{,}6$.

\section{Когда удобнее таблица исходов}
Если есть две игры и за каждую можно получить фиксированное число очков, компактная таблица часто нагляднее большого дерева. В каждой игре: поражение~--- $0$ очков, ничья~--- $1$, победа~--- $3$. Вероятность ничьей равна $0{,}1$, а победы и поражения~--- по $0{,}45$.

\begin{center}
\renewcommand{\arraystretch}{1.4}
\begin{tabular}{c|ccc}
\toprule
& \multicolumn{3}{c}{\textbf{Вторая игра}}\\
\textbf{Первая игра} & 0 $(0{,}45)$ & 1 $(0{,}1)$ & 3 $(0{,}45)$\\
\midrule
0 $(0{,}45)$ & \textbf{0} & \textbf{1} & \textbf{3}\\
1 $(0{,}1)$ & \textbf{1} & \textbf{2} & 4\\
3 $(0{,}45)$ & \textbf{3} & 4 & 6\\
\bottomrule
\end{tabular}
\end{center}
В клетках стоит \emph{сумма очков}. Для выхода в следующий круг нужно не менее $4$ очков. Неудачными исходами служат клетки с суммой $0$, $1$, $2$ или $3$. Игры считаются независимыми.
\[
\begin{aligned}
P(\text{не выйти})
&=0{,}45^2+2\cdot0{,}45\cdot0{,}1+0{,}1^2+2\cdot0{,}45^2\\
&=0{,}7075\approx\boxed{0{,}71}.
\end{aligned}
\]
Можно проверить результат через противоположное событие: выход возможен при сочетаниях «победа и ничья», «ничья и победа» или «две победы».

\section{Повторные попытки: первый успех}
Если одна попытка успешна с вероятностью $p$, а неуспешна с вероятностью $1-p$, то вероятность хотя бы одного успеха за $n$ \emph{независимых} попыток равна
\[
\boxed{P_n=1-(1-p)^n}.
\]
В частности, при вероятности попадания $0{,}8$ получаем
\[
P_1=0{,}8,\quad P_2=1-0{,}2^2=0{,}96,\quad
P_3=1-0{,}2^3=0{,}992.
\]
Следовательно, чтобы вероятность уничтожения цели была \emph{не менее} $0{,}99$, достаточно трёх выстрелов. После уничтожения цели стрельба прекращается. Формула учитывает вероятность успеха в течение первых $n$ разрешённых попыток.

\smalltopic{Если требуется минимальное количество попыток}
Нужно найти \emph{наименьшее натуральное} $n$, для которого $P_n$ достигает требуемого порога, и проверить, что при $n-1$ порог ещё не достигнут. Для порога «более $0{,}99$» используется строгое неравенство $P_n>0{,}99$.

\section{Извлечение предметов без возвращения}
В ящике $4$ красных и $2$ синих фломастера. Их извлекают последовательно, не возвращая обратно. Чтобы \emph{первый синий} оказался третьим, последовательность цветов должна быть строго такой:
\[
\text{красный}\;\to\;\text{красный}\;\to\;\text{синий}.
\]
После первого красного осталось $5$ фломастеров (из них $3$ красных), после второго~--- $4$ (из них $2$ синих). Поэтому
\[
P=\frac46\cdot\frac35\cdot\frac24=\boxed{\frac15}.
\]
\important{Типичная ошибка:} при извлечении без возвращения знаменатель меняется: $6$, затем $5$, затем $4$. Вероятности следующих извлечений зависят от того, что уже произошло.

\section{Короткий контроль перед ответом}
\begin{itemize}
\item Верно ли определено искомое событие: болезнь или результат анализа, успех одной попытки или всей серии?
\item Все ли пути учтены? На каждой развилке сумма вероятностей всех исходов равна единице?
\item Верно ли перемножены условные вероятности вдоль пути и сложены вероятности несовместимых путей?
\item Проверены ли формулировки «не менее» и «строго больше», округление и допустимость ответа от $0$ до $1$?
\end{itemize}

\clearpage
\thispagestyle{fancy}
\begin{center}
{\Large\bfseries\textcolor{darkaccent}{Алгоритм выбора способа решения}}\\[5mm]
\small \textcolor{muted}{Используйте после внимательного чтения полного условия}
\end{center}
\vspace{6mm}
\begin{center}
\begin{tikzpicture}[
>=Latex,
start/.style={ellipse,draw=accent,very thick,align=center,text width=37mm,minimum height=12mm,font=\small},
proc/.style={rectangle,rounded corners=2pt,draw=accent,thick,align=center,text width=40mm,minimum height=16mm,inner sep=5pt,font=\small},
dec/.style={diamond,draw=accent,thick,aspect=2.3,align=center,text width=49mm,inner sep=2pt,font=\small},
arr/.style={-Latex,thick,draw=darkaccent}]
\node[start] (beg) at (0,0) {Прочитайте условие};
\node[proc,text width=59mm] (events) at (0,-1.8) {Выделите этапы опыта, исходы и искомое событие};
\node[dec] (question) at (0,-5.0) {Важен порядок этапов?};
\node[proc,text width=44mm] (tree) at (-3.2,-8.0) {\textbf{Дерево}\\[1mm]Подпишите этапы и вероятности на ветвях};
\node[proc,text width=44mm] (table) at (3.2,-8.0) {\textbf{Таблица}\\[1mm]Составьте таблицу исходов и очков};
\node[proc,text width=99mm,minimum height=19mm] (calc) at (0,-11.1) {Отметьте подходящие исходы. Для каждого пути умножьте вероятности его этапов; затем сложите вероятности разных путей};
\node[start] (end) at (0,-13.4) {Проверьте результат};
\draw[arr] (beg)--(events);
\draw[arr] (events)--(question);
\draw[arr] (question.west) -| node[pos=.25,above,font=\small]{да} (tree.north);
\draw[arr] (question.east) -| node[pos=.25,above,font=\small]{нет} (table.north);
\draw[arr] (tree.south) -- ([xshift=-25mm]calc.north);
\draw[arr] (table.south) -- ([xshift=25mm]calc.north);
\draw[arr] (calc)--(end);
\end{tikzpicture}
\end{center}
\clearpage
\section*{Тренировочный блок}
\textcolor{muted}{Путь А: 10 задач. Решите последовательно, начиная с простых. Условия приведены полностью. В ответе запишите вероятность или искомую долю; при необходимости покажите дерево, таблицу или уравнение.}

\subsection*{Базовый уровень}
\begin{enumerate}[label=\textbf{\arabic*.},wide=0pt]
\Needspace{10\baselineskip}\item \textbf{Погода в Волшебной стране.} В Волшебной стране в течение каждого дня погода бывает только двух типов: хорошая или замечательная. Известно, что с вероятностью $0{,}8$ погода на следующий день сохраняет прежний тип, а с вероятностью $0{,}2$ меняется на другой. Вероятности переходов одинаковы для обоих типов погоды и не зависят от предшествующих дней, если известна погода текущего дня. Сегодня, $3$ июля, погода замечательная. Найдите вероятность того, что $5$ июля погода также будет замечательной.

\Needspace{10\baselineskip}\item \textbf{Прогулка к озеру.} Турист подходит к развилке, на которой наугад выбирает одну из двух дорожек: левую или правую. Обе дорожки выбираются с одинаковой вероятностью. Если турист пошёл налево, то на следующей развилке перед ним будет четыре равновероятных продолжения, причём три ведут к озеру, а одно~--- к парковке. Если турист пошёл направо, то на следующей развилке перед ним будет три равновероятных продолжения, причём одно ведёт к озеру, а два~--- к парковке. Турист всегда выбирает на второй развилке случайную дорожку и после второго выбора доходит до конечного пункта. Найдите вероятность того, что турист окажется у озера.

\Needspace{10\baselineskip}\item \textbf{Повторное исследование.} В лаборатории исследуют образец, про который заранее известно, что он не содержит опасного вещества. Каждое исследование может ошибочно дать положительный результат с вероятностью $0{,}02$ и отрицательный с вероятностью $0{,}98$. Если первое исследование оказалось отрицательным, образец исследуют повторно. После положительного результата дальнейшее исследование прекращают. Считайте, что при условии отсутствия вещества результаты двух исследований независимы. Найдите вероятность того, что хотя бы одно проведённое исследование окажется положительным.

\Needspace{12\baselineskip}\subsection*{Средний уровень}
\Needspace{10\baselineskip}\item \textbf{Положительный медицинский тест.} В клинику поступают пациенты с подозрением на заболевание. По статистике действительно больны $8\%$ поступающих пациентов, остальные здоровы. У действительно больного пациента тест даёт положительный результат с вероятностью $0{,}85$, а у здорового пациента ложный положительный результат возникает с вероятностью $0{,}03$. Выбирают одного пациента из числа поступивших в клинику и проводят ему один тест. Найдите вероятность того, что результат этого теста будет положительным.

\Needspace{10\baselineskip}\item \textbf{Неизвестная доля больных.} Всем пациентам диагностического центра проводят один и тот же анализ. У человека, действительно имеющего заболевание, результат анализа положителен с вероятностью $0{,}82$. У здорового человека положительный результат может возникнуть ошибочно с вероятностью $0{,}02$. Известно, что среди всех пациентов центра положительный результат анализа получают $7{,}5\%$ человек. Найдите долю пациентов, которые действительно имеют заболевание. Считайте, что указанные условные вероятности одинаковы для всех соответствующих пациентов.

\Needspace{10\baselineskip}\item \textbf{Яйца из двух хозяйств.} Торговая компания закупает яйца у двух птицефабрик. В первой птицефабрике к высшей категории относятся $60\%$ яиц, во второй~--- $25\%$. Известно, что среди всех яиц, приобретённых компанией у этих двух птицефабрик, доля яиц высшей категории равна $39\%$. Каждое яйцо поступило ровно с одной из двух птицефабрик. Найдите вероятность того, что случайно выбранное среди закупленных компанией яиц яйцо поступило с первой птицефабрики.

\Needspace{12\baselineskip}\subsection*{Повышенный уровень}
\Needspace{10\baselineskip}\item \textbf{Две игры футбольной команды.} Чтобы пройти в следующий круг турнира, футбольной команде необходимо набрать не менее $4$ очков по итогам двух игр. За победу команда получает $3$ очка, за ничью~--- $1$ очко, за поражение~--- $0$ очков. В каждой игре вероятность ничьей равна $0{,}2$, а вероятности победы и поражения одинаковы. Результаты двух игр независимы друг от друга. Найдите вероятность того, что команда \emph{не сможет} выйти в следующий круг соревнований. В ответе укажите десятичную дробь.

\Needspace{10\baselineskip}\item \textbf{Автоматическая система стрельбы.} Автоматическая система делает выстрелы по неподвижной цели. После первого уничтожения цели система немедленно прекращает стрельбу. Если выстрел не уничтожил цель, выполняется следующий выстрел. Каждый выстрел уничтожает ещё не поражённую цель с вероятностью $0{,}65$ независимо от результатов предыдущих выстрелов. Какое \emph{наименьшее число выстрелов} следует разрешить системе, чтобы вероятность уничтожения цели за это число выстрелов была не менее $0{,}98$?

\Needspace{12\baselineskip}\subsection*{Сложный уровень}
\Needspace{10\baselineskip}\item \textbf{Фломастеры без возвращения.} В непрозрачном ящике лежат пять красных и три синих фломастера, одинаковых на ощупь. Фломастеры последовательно извлекают из ящика наугад, не рассматривая содержимое ящика и не возвращая уже извлечённые фломастеры. Найдите вероятность того, что \emph{первый синий фломастер} будет извлечён именно третьим по счёту. Требуется учесть изменение количества красных и синих фломастеров после каждого извлечения.

\Needspace{10\baselineskip}\item \textbf{Пенсионеры среди жителей города.} В городе $45\%$ жителей составляют мужчины, а остальные $55\%$~--- женщины. По данным городского статистического обследования пенсионеры составляют $13{,}75\%$ всего взрослого населения, а среди женщин доля пенсионерок равна $16\%$. Считайте, что приведённые доли мужчин и женщин относятся к тому же взрослому населению, которое охватывает обследование. Для социологического опроса случайно выбирают одного мужчину из числа взрослых мужчин города. Найдите вероятность того, что выбранный мужчина является пенсионером.
\end{enumerate}

\clearpage
\section*{Ответы к тренировочному блоку}
\textcolor{muted}{Сначала выполните задания самостоятельно. Таблица предназначена для самопроверки; развёрнутые решения можно восстановить по правилам и примерам чек-листа.}
\vspace{6mm}
\begin{center}
\renewcommand{\arraystretch}{1.65}
\begin{tabularx}{\linewidth}{@{}>{\bfseries}p{13mm}X>{\raggedleft\arraybackslash}p{42mm}@{}}
\toprule
№ & Краткий ориентир & Ответ\\
\midrule
1 & Погода через два дня & $0{,}68$\\
2 & Два направления маршрута & $\dfrac{13}{24}$\\
3 & Положительный результат с повтором & $0{,}0396$\\
4 & Общая вероятность положительного теста & $0{,}0956$\\
5 & Доля действительно больных & $0{,}06875$ ($6{,}875\%$)\\
6 & Доля продукции первой птицефабрики & $0{,}4$ ($40\%$)\\
7 & Неудачный итог двух игр & $0{,}68$\\
8 & Минимальное число выстрелов & $4$\\
9 & Первый синий фломастер~--- третий & $\dfrac{5}{28}$\\
10 & Доля пенсионеров среди мужчин & $0{,}11$ ($11\%$)\\
\bottomrule
\end{tabularx}
\end{center}
\vspace{10mm}
\smalltopic{Финальная памятка}
\begin{itemize}
\item \textbf{«И затем»} в одном пути дерева: умножьте условные вероятности.
\item \textbf{«Или»} для непересекающихся путей: сложите их вероятности.
\item \textbf{Дана общая частота события}: при неизвестной доле группы введите $x$ и составьте уравнение полной вероятности.
\item \textbf{Результат анализа}: разделяйте состояние пациента и вывод прибора.
\item \textbf{Несколько игр}: перечислите пары исходов в таблице и выделите нужную сумму очков.
\item \textbf{Повторные попытки}: удобнее посчитать вероятность отсутствия успеха во всех попытках.
\item \textbf{Выбор без возвращения}: после каждого извлечения меняйте количество оставшихся предметов.
\end{itemize}
\end{document}
